% === Begin Label Data ===
% ID: 787
% 难度星级: 
% 标签: 
% 备注: 
% 组卷引用次数: 0
% === End  Label Data ===

\begin{problem}{2015}{G}{上海卷（文）}{13}{向量}
已知平面向量 $ \vec{a} $, $ \vec{b} $, $ \vec{c} $ 满足 $ \vec{a}\perp\vec{b} $, 且 $ \{|\vec{a}|,|\vec{b}|,|\vec{c}|\}=\{1,2,3\} $, 则 $ |\vec{a}+\vec{b}+\vec{c}| $ 的最大值是\underline{\hspace{4em}}.
\end{problem}

\begin{answer}
$ 3+\sqrt{5} $
\end{answer}

\begin{solutions}
思路 1：$ |\vec{a}+\vec{b}+\vec{c}|=\sqrt{(\vec{a}+\vec{b}+\vec{c})^{2}}=\sqrt{|\vec{a}|^{2}+|\vec{b}|^{2}+|\vec{c}|^{2}+2\vec{a}\cdot\vec{b}+2\vec{b}\cdot\vec{c}+2\vec{c}\cdot\vec{a}} $.  
由于 $ \vec{a}\perp\vec{b} $, 且 $ \{|\vec{a}|,|\vec{b}|,|\vec{c}|\}=\{1,2,3\} $, 所以 $ \vec{a}\cdot\vec{b}=0 $, $ |\vec{a}|^{2}+|\vec{b}|^{2}+|\vec{c}|^{2}=1+4+9=14 $,  
$ \vec{b}\cdot\vec{c}+\vec{c}\cdot\vec{a}=(\vec{a}+\vec{b})\cdot\vec{c}\leqslant|\vec{a}+\vec{b}||\vec{c}| $, 当 $ \vec{a}+\vec{b} $ 与 $ \vec{c} $ 同向时取等号, 而 $ |\vec{a}+\vec{b}|=\sqrt{|\vec{a}|^{2}+|\vec{b}|^{2}} $, 经过奇数可知 $ \{|\vec{a}|,|\vec{b}|\}=\{1,2\} $, 且 $ |\vec{c}|=3 $ 时, $ |\vec{a}+\vec{b}||\vec{c}| $ 有最大值 $ 3\sqrt{5} $. 因此, 所求的最大值为 $ \sqrt{14+2\times3\sqrt{5}} $, 即 $ 3+\sqrt{5} $.  

思路 2：设 $ \vec{d}=\vec{a}+\vec{b} $, 则 $ |\vec{d}|^{2}=|\vec{a}|^{2}+|\vec{b}|^{2}+2\vec{a}\cdot\vec{b}=|\vec{a}|^{2}+|\vec{b}|^{2} $.  
于是 $ |\vec{a}+\vec{b}+\vec{c}|=|\vec{d}+\vec{c}|\leqslant|\vec{d}|+|\vec{c}|=\sqrt{|\vec{a}|^{2}+|\vec{b}|^{2}}+|\vec{c}| $,  
等号成立当且仅当 $ \vec{d} $ 与 $ \vec{c} $ 共线. 逐一代入可知当 $ |\vec{a}|=1 $, $ |\vec{b}|=2 $, $ |\vec{c}|=3 $ 时, 欲求表达式取最大值 $ 3+\sqrt{5} $.
\end{solutions}